\clearpage
\SourceLeaf{374}%
\begin{center}{\large [DIE OBERFLÄCHE DES ELLIPSOIDS.]\par}\end{center}
\begin{center}[1.]\end{center}

Wenn jeder Punkt einer Fläche durch denjenigen Punkt einer Kugelfläche mit Halbmesser 1 repräsentirt wird, dessen Radius mit der Normale des ersten Punkts parallel ist, und ein Element der Kugelfläche $d\sigma$ einem Element der Fläche $ds$ correspondirt, so ist
\[
\frac{d\sigma}{ds}=\frac{\dfrac{ddz}{dx^2}\dfrac{ddz}{dy^2}-\left(\dfrac{ddz}{dx\,dy}\right)^2}{\left(1+\left(\dfrac{dz}{dx}\right)^2+\left(\dfrac{dz}{dy}\right)^2\right)^2}=\frac1{rr'},
\]
wenn $r$ und $r'$ die beiden äussersten Krümmungshalbmesser sind.
\begin{center}[2.]\end{center}

Wenn $d\sigma$ den Raum auf der Himmelskugel bedeutet, welcher dem Elemente $ds$ der Oberfläche des Ellipsoids
\[
\left[\frac{xx}{aa}+\frac{yy}{bb}+\frac{zz}{cc}=1\right]
\]
entspricht, so wird
\[
ds=aabbcc\left(\frac{xx}{a^4}+\frac{yy}{b^4}+\frac{zz}{c^4}\right)^2d\sigma.
\]
Sind hier $\xi,\eta,\zeta$ die Coordinaten des Punktes der Himmelskugel, so wird
